% =====================================================
% 02_related.tex — short related work, ~0.75 page
% =====================================================

\paragraph{GAN interpretation and attribute editing.}
\textbf{Unit-level / segmentation-based dissection}
\cite{Bau2019GANDissection} probes individual generator units via
attached segmentation classifiers. \textbf{Boundary-based editing}
\cite{Shen2020InterFaceGAN, Yang2020HiGAN} learns linear hyperplanes
in $W^+$ for semantic attributes; we use the resulting boundaries as
inputs to our analysis. \textbf{Unsupervised discovery} approaches:
PCA on activations~\cite{Harkonen2020GANSpace}, closed-form
factorisation of the affine layer~\cite{Shen2021SeFa}, channel-wise
style manipulation~\cite{Wu2021StyleSpace}, contrastive direction
learning~\cite{Yuksel2021LatentCLR, Ren2022DisCo}, and
deformation-field editing~\cite{Pan2023DragGAN}. We compare against
all of these on attribute rediscovery and editing-task head-to-head
(§\ref{sec:experiments}).

\paragraph{Diffusion-model interpretation.}
The diffusion analogue of $W$-space is the U-Net bottleneck h-space
introduced by Asyrp~\cite{Kwon2023Asyrp}, with homogeneity, linearity,
and timestep-consistency reported by the authors. Park et
al.~\cite{Park2023Riemannian} constructs a first-order pullback
Riemannian metric on h-space via Jacobian SVD and uses it for
editing; \textbf{they consider only the metric, not the curvature}.
Prompt-side editing methods --- Prompt-to-Prompt~\cite{Hertz2023P2P},
null-text inversion~\cite{Mokady2023NullText},
SEGA~\cite{Brack2023SEGA} --- operate in attention space rather than
on the generator's image manifold. DAAM~\cite{Tang2023DAAM} provides
cross-attention saliency at the prompt level. We use h-space (Asyrp)
as the W-analog, compute the second-fundamental form there, and
compare our gradient saliency to DAAM directly.

\paragraph{Score-network geometry of diffusion.}
Stanczuk et al.~\cite{Stanczuk2024IntrinsicDim} measure rank
deficiency of the score Jacobian to recover the intrinsic data
dimension. Higher-order ODE solvers
(GENIE~\cite{Dockhorn2022GENIE}) use JVPs through the U-Net for
better numerical accuracy. Our use of forward-mode AD through the
DDIM sampling chain extends the GENIE recipe to second-order
geometric measurement of the rendered image, complementing the
score-side analyses of~\cite{Stanczuk2024IntrinsicDim}.

\paragraph{Classical saliency, image attribution, and Grad-CAM.}
CAM~\cite{Zhou2016CAM} and Grad-CAM~\cite{Selvaraju2017GradCAM}
attribute classifier decisions to spatial regions via either pooled
weights or class-score gradients. Our pushforward saliency parallels
Grad-CAM in spirit but operates on the generator alone, with no
classifier --- the source of curvature is the model itself, not a
recognition head. As §\ref{sec:experiments} shows, CLIP-Grad-CAM
saliency and our JVP saliency are cross-domain near-orthogonal
($|r| < 0.11$ on bedroom, $< 0.02$ on FFHQ), establishing them as
complementary rather than competing measurements.

\paragraph{Differential geometry of neural-network manifolds.}
Natural gradient~\cite{Amari1998NaturalGradient} and geometric deep
learning~\cite{Bronstein2021GeometricDL} formalise the ambient
geometry of \emph{parameter} space and \emph{input} domains. Our
setting is different: we study the geometry of the manifold of
\emph{outputs} learned by a fixed generator. Zhu et
al.~\cite{Zhu2016ManipulationOnManifold} introduced the "natural
image manifold" formulation; we provide a tractable second-order
estimator on this manifold and connect curvature to disentanglement
and compositional editing.

\paragraph{Forward-mode autodifferentiation.}
The Pearlmutter trick~\cite{Pearlmutter1994HVP} for Hessian-vector
products underlies modern second-order optimisation
\cite{Martens2015KFAC}; PyTorch \texttt{torch.func.jvp} and
\texttt{jax.jvp} provide composable forward AD. We compose JVP twice
to obtain the per-pixel second derivative in one pass (§4), which is
$10^3$-$10^4 \times$ faster than per-pixel reverse mode for the
high-dimensional-output regime studied here.

\paragraph{Encoders, inversion, and chart consistency.}
pSp~\cite{Richardson2021pSp}, e4e~\cite{Tov2021e4e}, and the
survey~\cite{Xia2022InversionSurvey} establish reconstruction-error
as the de facto encoder-quality metric. We argue (§5.5) that
\emph{first-order} chart consistency --- agreement of saliency
derivatives at the encoded code with those at the ground-truth code
--- is a strictly stronger signal: on bedroom, the saliency
correlation rises 45$\times$ while recon MSE only falls 17\% across
40k iterations of synthetic-supervision training.

A full survey of $\sim$60 additional references --- spanning
StyleGAN-NADA, EditGAN~\cite{Ling2021EditGAN}, NoiseCLR, LELSD,
WarpedGANSpace, Concept Algebra, LEDITS, and Hessian-Penalty --- is
in the supplementary material.
